



\iffalse
\subsection{Datasets and Validation Strategy}

Our model training and evaluation strategy follows a robust cross-dataset approach to ensure generalizability across different recording conditions and gesture contexts. We employ a carefully structured dataset split:

\subsubsection{Cross-Dataset Validation}
To rigorously assess our model's generalization capabilities, we employ a two-stage validation strategy:

\paragraph{Held-out Validation} We reserve 10\% of the SAGA corpus (approximately 2.5 dialogues) for validation during training. This dataset was chosen for validation because it provides the most detailed gesture annotations and multi-modal recordings, allowing for precise evaluation of model performance.

\paragraph{Cross-Dataset Testing} We use the MULTISIMO corpus as our primary test set, specifically because:
\begin{itemize}
   \item Its collaborative task-based setting differs significantly from our training contexts
   \item It features multiple participants and varying camera angles
   \item It contains naturalistic gestures that challenge the model's ability to generalize
\end{itemize}

This cross-dataset evaluation approach helps ensure that our model performs robustly across different gesture contexts, recording conditions, and speaker styles, rather than overfitting to any particular dataset's characteristics.
\fi
\subsection{Model and Training}
\begin{itemize}
    \item Framework: TensorFlow.
    \item Skeleton tracking using Mediapipe Holistic.
    \item Training metrics: Accuracy, validation accuracy, etc.
    % : Which papers detail Mediapipe ,open pose, ....TensorFlow training methodologies?
\end{itemize}
\iffalse


\subsection{Computational Efficiency}
The model demonstrates strong computational efficiency, making it suitable for real-world applications on consumer hardware. Running on a standard Dell LAtituted (64gb RAM notebook /.....%CPU (Intel i7-10700K), the system processes video input at an average rate of ....%25 frames per second, with the following component-wise latencies:


\subsubsection{Pose features: Feature extractions}
We prepare 29 features from the pose data, which are essentially normalized pose-embeddings. 

The model was trained using the following configuration:
\begin{itemize}
    \item Batch size: 32
    \item Learning rate: 0.001 with Adam optimizer
    \item Training epochs: 100 with early stopping
    \item Input sequence length: 32 frames
    \item Loss function: Categorical cross-entropy
\end{itemize}

We monitor several metrics during training:
\begin{itemize}
    \item Frame-level accuracy
    \item Gesture-level F1 score
    \item Temporal intersection over union (tIoU)
    \item Confusion matrix across gesture types
\end{itemize}

To prevent overfitting, we employ regularization techniques including dropout (rate = 0.5) and L2 regularization (λ = 0.01). Data augmentation techniques such as random temporal scaling and random noise addition to landmark positions are applied during training to improve model robustness.
\fi